\chapter{Vergleich}
\label{vergleich}

\section{Evaluierung}
Der finale Teil der Arbeit beschäftigt sich mit dem Vergleich der Performance,
des Systems. Das System wurde mit vier verschiedenen LLM Modellen implementiert und 
wird im ersten Schritt mit dem Framework RAGAS (siehe \ref{s:intro:ragas}) miteinander 
verglichen. Nach diesem Schritt werden die zwei bestabschneidensten Modelle 
für die finale Bewertung von Mitarbeitern aus der Techdivision miteinander 
verglichen. 

\section{RAGAS Vergleich}
Die erste Ausscheidungsrunde wird über das Benchmarking Frameworks entschieden.
Die vier LLM Modelle die hier zur Auswahl stehen sind finetuned von Metas Llama3,
Mistrals Mistral 2 und Googles Gemma. Dazu kommt noch die originale Version von 
Microsofts Phi 3.\par

\subsection{Testset}
Der erste Schritt für den Vergleich ist es Testsets herzustellen anhand
derer Ragas den Benchmark Test durchführen kann.\par
Die Variante die hierfür verwendet wird ist es einen TestsetGenerator per 
\par\verb+TestsetGenerator.from_langchain()+
herzustellen und mit diesem dann das Testset erstellen zu lassen.\par
Als Vorbereitung hierfür werden Daten benötigt. 
Die Daten die als Dokument geladen werden haben alle ein 
Metadata Dictionary. Damit der Generator erkennen kann,
welche Daten zusammengehören brauchen die Dokumente in 
diesem Metadata Dictionary ein Key Value Pair mit dem Namen
''filename``. Der Wert des ''filename`` Schlüssels wird verwendet 
um bei der Generation der Testdaten zusammengehörige 
Dokumententeile zusammenzufügen.\par
Nach dem Vorbereiten der Daten kann man damit beginnen den 
Generator zu erstellen.
Für diesen Generator werden drei Parameter verwendet.
Die \verb+generator_llm+, \verb+critic_llm+ und \verb+embeddings+.
Als Embedding Model benutzen wir ebenfalls das Model 
"intfloat/multilingual-e5-large", wie für den Chatbot und 
den Importer. 
Als Critic und Generator LLM werden wir zwei verschiedene LLMs
nutzen. Da die Geschwindigkeit der Erstellung dieser Daten
keine große Rolle spielt werden hierfür große Modelle genutzt.
Das Generator Modell wird eine auf Deutsch finetuned Version
von Mistrals Mixtral 8x7B sein und zwar 
"VAGOsolutions/SauerkrautLM-Mixtral-8x7B" \cite{VAGOsolutionsSauerkrautLMMixtral8x7BHugging2023}.
Als Critic Modell wird eine Version von Microsofts Phi 3 genutzt, welche
ebenfalls auf Deutsch finetuned wurde. Die Version die genutzt wird ist 
"VAGOsolutions/SauerkrautLM-Phi-3-medium", welche eine finetuned Version 
von dem Modell "unsloth/Phi-3-medium-4k-instruct" ist. \par
Nach dem Erstellen des Generators und dem vor und aufbereiten der 
Daten kann man das Testset generieren lassen über die Funktion 
\verb+generate_with_langchain_docs()+. 
Für die Generierung kann man drei Parameter setzen.
Als erstes die \verb+documents+ von denen man die Testsets 
generieren will.
Hier werden Documents verwendet die über den \verb+ConfluenceLoader+
von Langchain geladen werden. Parallel wie beim Importer (\ref{s:dev:imp}).
Dazu kann man bestimmen, wie viele Elemente im Set enthalten sein sollen.
Diese \verb+test_size+ bestimmt wie viele Fragen vom \verb+TestsetGenerator+ generiert werden.
Dazu kann man die Verteilung der Fragen bestimmen.
Die Fragen die durch den Generator erstellt werden können dazu drei verschiedene
Typen haben. Diese drei Typen sind: \verb+simple+, \verb+reasoning+ und \verb+multi_context+.
Diese kann man der \verb+generate+ Funktion als dict mitgeben mit der Verteilung 
die man verwenden möchte. Hier werden diese Werte verwendet: \verb+simple: 0.5+,
\verb+reasoning: 0.2+ und \verb+multi_context: 0.3+.\par
Das Testset welches hierbei generiert wird kann man mittels 
\verb+to_pandas()+ als pandas Dataframe ausgeben.
Wenn man dieses Dataframe betrachtet sieht man, 4 relevante Spalten.\par
\verb+question+: Die Frage die dem Chatbot für das Benchmarking gestellt wird. \par
\verb+context+: Der Context aus dem die Frage erstellt wurde. \par
\verb+answer+: Die Antwort auf die Frage. \par
\verb+question_type+: Der Question Typ, dessen Verhältnis in der \verb+generate+ Funktion bestimmt wurde.\par
Mit diesem Testset kann man nun die Chain bewerten.

\subsection{Evaluation}
Nachdem das Testset generiert wurde kann es zum bewerten
übergehen. Für die Bewertung gibt es verschiedene Metriken, die man 
betrachten kann. 

\subsubsection{Metriken}
\begin{itemize}
    \item Faithfulness
    \item Answer relevancy
    \item Answer semantic similarity
    \item Answer correctness
    \item Context recall
    \item Context precision
    \item Context relevancy
    \item Context entity recall
\end{itemize}

\section{Umfrage Vergleich}
Hierfür wurden 5 Mitarbeiter aus dem Unternehmen mit einem Fragebogen 
ausgestattet und ihnen wurde Zugang zum System gegeben. 

\subsection{Fragebogen}
Der Fragebogen wurde so erstellt, dass es keine genauen Fragen gibt, 
sondern nur Themengebiete, die abgefragt werden sollen. Das wichtige 
beim Auswählen der Themengebiete war, dass es Informationen waren, die 
auf der Confluence Instanz, auf die der Chatbot Zugriff hatte, verfügbar 
waren.